% GENERATED FILE. Edit canonical lessons, then run npm run generate:books.
% canonical-source-hash: fnv1a64:be1d69cb3a8898c0
% canonical-lessons: ML-C329-lekhanam, ML-C329-kavita, ML-C329-natakam, ML-C329-phutbol, ML-C329-krikkarru

\chapter{Article, Poem, Play, Football, Cricket}
\label{ch:ml-article-poem-play-football-cricket}
\hlchaptermodality{329}

\begin{chapteropening}
\textbf{By the end of this chapter:} \emph{I can say an article, a poem, a play, football and cricket.}

Complete the last lesson of chapter 329: I can say an article, a poem, a play, football and cricket.
\end{chapteropening}

\section[\texorpdfstring{\ml{ലേഖനം} (lēkhanaṁ)}{ലേഖനം (lēkhanaṁ)}]{\ml{ലേഖനം} (lēkhanaṁ) — an article}

\label{lesson:ML-C329-lekhanam}



\begin{quote}
Before the new one: say the Malayalam for an axe, then the Malayalam for a spade.
\end{quote}

\subsection*{You'll want to know: \ml{ലേഖനം}}
\textbf{\ml{ലേഖനം}} — \emph{lēkhanaṁ} — \textquotedblleft{}an article\textquotedblright{}.

\begin{cousinweb}
\textbf{\ml{ലേഖനം}} (\emph{lēkhanaṁ}) $\leftarrow$ Sanskrit \emph{lēkhana}, \textquotedblleft{}scraping; writing\textquotedblright{}, taken over whole as a tatsama, with Malayalam's noun ending \emph{-ṁ} (CDIAL 11103). Writing began as scratching lines.
\end{cousinweb}

\begin{grammarlens}[title={What you've built}]
One more word for this chapter.
\end{grammarlens}

\subsection*{Your turn}
\begin{itemize}
\raggedright
  \item \emph{Say it:} \emph{lēkhanaṁ}
  \item \emph{Say it:} \emph{lēkhanaṁ}, once more
  \item \emph{Say it:} say \emph{lēkhanaṁ}
  \item \emph{Recall:} say the Malayalam for an axe, then the Malayalam for a spade, then say \emph{lēkhanaṁ} again
\end{itemize}

\begin{tcolorbox}[breakable,colback=teal!4,colframe=teal!35!black,title={Before you move on}]
What does \ml{ലേഖനം} mean? (\textquotedblleft{}An article\textquotedblright{}.) Say it once more.
\end{tcolorbox}
\section[\texorpdfstring{\ml{കവിത} (kavita)}{കവിത (kavita)}]{\ml{കവിത} (kavita) — a poem}

\label{lesson:ML-C329-kavita}



\begin{quote}
Before the new one: say the Malayalam for a spade, then the Malayalam for an article.
\end{quote}

\subsection*{You'll want to know: \ml{കവിത}}
\textbf{\ml{കവിത}} — \emph{kavita} — \textquotedblleft{}a poem\textquotedblright{}.

\begin{cousinweb}
\textbf{\ml{കവിത}} (\emph{kavita}) $\leftarrow$ Sanskrit \emph{kavitā}, \textquotedblleft{}poetry\textquotedblright{}, taken over whole as a tatsama. It is built on \emph{kavi}, \textquotedblleft{}poet\textquotedblright{}.
\end{cousinweb}

\begin{grammarlens}[title={What you've built}]
One more word for this chapter.
\end{grammarlens}

\subsection*{Your turn}
\begin{itemize}
\raggedright
  \item \emph{Say it:} \emph{kavita}
  \item \emph{Say it:} \emph{kavita}, once more
  \item \emph{Say it:} say \emph{kavita}
  \item \emph{Recall:} say the Malayalam for a spade, then the Malayalam for an article, then say \emph{kavita} again
\end{itemize}

\begin{tcolorbox}[breakable,colback=teal!4,colframe=teal!35!black,title={Before you move on}]
What does \ml{കവിത} mean? (\textquotedblleft{}A poem\textquotedblright{}.) Say it once more.
\end{tcolorbox}
\section[\texorpdfstring{\ml{നാടകം} (nāṭakaṁ)}{നാടകം (nāṭakaṁ)}]{\ml{നാടകം} (nāṭakaṁ) — a play, a drama}

\label{lesson:ML-C329-natakam}



\begin{quote}
Before the new one: say the Malayalam for an article, then the Malayalam for a poem.
\end{quote}

\subsection*{You'll want to know: \ml{നാടകം}}
\textbf{\ml{നാടകം}} — \emph{nāṭakaṁ} — \textquotedblleft{}a play, a drama\textquotedblright{}.

\begin{cousinweb}
\textbf{\ml{നാടകം}} (\emph{nāṭakaṁ}) $\leftarrow$ Sanskrit \emph{nāṭaka}, \textquotedblleft{}dancer, actor; a play\textquotedblright{}, taken over whole as a tatsama, with Malayalam's noun ending \emph{-ṁ} (CDIAL 7045). It is built on \emph{naṭa}, \textquotedblleft{}actor, dancer\textquotedblright{}.
\end{cousinweb}

\begin{grammarlens}[title={What you've built}]
One more word for this chapter.
\end{grammarlens}

\subsection*{Your turn}
\begin{itemize}
\raggedright
  \item \emph{Say it:} \emph{nāṭakaṁ}
  \item \emph{Say it:} \emph{nāṭakaṁ}, once more
  \item \emph{Say it:} say \emph{nāṭakaṁ}
  \item \emph{Recall:} say the Malayalam for an article, then the Malayalam for a poem, then say \emph{nāṭakaṁ} again
\end{itemize}

\begin{tcolorbox}[breakable,colback=teal!4,colframe=teal!35!black,title={Before you move on}]
What does \ml{നാടകം} mean? (\textquotedblleft{}A play, a drama\textquotedblright{}.) Say it once more.
\end{tcolorbox}
\section[\texorpdfstring{\ml{ഫുട്ബോൾ} (phuṭbōḷ)}{ഫുട്ബോൾ (phuṭbōḷ)}]{\ml{ഫുട്ബോൾ} (phuṭbōḷ) — football}

\label{lesson:ML-C329-phutbol}



\begin{quote}
Before the new one: say the Malayalam for a poem, then the Malayalam for a play.
\end{quote}

\subsection*{You'll want to know: \ml{ഫുട്ബോൾ}}
\textbf{\ml{ഫുട്ബോൾ}} — \emph{phuṭbōḷ} — \textquotedblleft{}football\textquotedblright{}.

\begin{cousinweb}
\textbf{\ml{ഫുട്ബോൾ}} (\emph{phuṭbōḷ}) $\leftarrow$ English \emph{football}.
\end{cousinweb}

\begin{grammarlens}[title={What you've built}]
One more word for this chapter.
\end{grammarlens}

\subsection*{Your turn}
\begin{itemize}
\raggedright
  \item \emph{Say it:} \emph{phuṭbōḷ}
  \item \emph{Say it:} \emph{phuṭbōḷ}, once more
  \item \emph{Say it:} say \emph{phuṭbōḷ}
  \item \emph{Recall:} say the Malayalam for a poem, then the Malayalam for a play, then say \emph{phuṭbōḷ} again
\end{itemize}

\begin{tcolorbox}[breakable,colback=teal!4,colframe=teal!35!black,title={Before you move on}]
What does \ml{ഫുട്ബോൾ} mean? (\textquotedblleft{}Football\textquotedblright{}.) Say it once more.
\end{tcolorbox}
\section[\texorpdfstring{\ml{ക്രിക്കറ്റ്} (krikkaṟṟŭ)}{ക്രിക്കറ്റ് (krikkaṟṟŭ)}]{\ml{ക്രിക്കറ്റ്} (krikkaṟṟŭ) — cricket}

\label{lesson:ML-C329-krikkarru}



\begin{quote}
Before the new one: say the Malayalam for a play, then the Malayalam for football.
\end{quote}

\subsection*{You'll want to know: \ml{ക്രിക്കറ്റ്}}
\textbf{\ml{ക്രിക്കറ്റ്}} — \emph{krikkaṟṟŭ} — \textquotedblleft{}cricket\textquotedblright{}.

\begin{cousinweb}
\textbf{\ml{ക്രിക്കറ്റ്}} (\emph{krikkaṟṟŭ}) $\leftarrow$ English \emph{cricket}.
\end{cousinweb}

\begin{grammarlens}[title={What you've built}]
One more word for this chapter.
\end{grammarlens}

\subsection*{Your turn}
\begin{itemize}
\raggedright
  \item \emph{Say it:} \emph{krikkaṟṟŭ}
  \item \emph{Say it:} \emph{krikkaṟṟŭ}, once more
  \item \emph{Say it:} say \emph{krikkaṟṟŭ}
  \item \emph{Recall:} say the Malayalam for a play, then the Malayalam for football, then say \emph{krikkaṟṟŭ} again
\end{itemize}

\begin{tcolorbox}[breakable,colback=teal!4,colframe=teal!35!black,title={Before you move on}]
What does \ml{ക്രിക്കറ്റ്} mean? (\textquotedblleft{}Cricket\textquotedblright{}.) Say it once more.
\end{tcolorbox}
